% ==========================================================================
%  Chapter 87 -- Thesis/Library Alignment and Stage-1 Research Roadmap
% ==========================================================================

\chapter{Thesis/Library Alignment and Stage-1 Research Roadmap}
\label{ch:thesis-library-stage1}

\paragraph{Normative status.}
This chapter records the first joint roadmap for the thesis and the library.
It does not replace the current normative multiscale contracts
(Chapter~\ref{ch:multiscale-publication-audit}); instead, it states how the
next architectural iteration must be executed so that the thesis narrative and
the implementation evolve consistently.

\section{Why a joint roadmap is necessary}

The current codebase already implements several ideas that the thesis only
states at a conceptual level: checkpointable solvers, typed multiscale
aggregation, explicit coupling sites, section-law diagnostics, and a material
factory boundary separating the multiscale core from Ko--Bathe/Menegotto
defaults.  Conversely, the thesis already develops a much stronger ontology for
continuum mechanics than the one currently enforced by the geometry/topology
layer in code.

This creates a productive but potentially dangerous asymmetry:
\begin{itemize}
  \item the code may overfit to one current realization of the method;
  \item the thesis may overstate capabilities that are still partial or
  placeholder; and
  \item future work on alternative local models (XFEM, DG, DPG, surrogate
  operators) may remain architecturally blocked if the current
  beam-section-focused contract becomes the de facto permanent interface.
\end{itemize}

The purpose of the present chapter is therefore not to promise a finished
architecture, but to define a disciplined route from the current
FE\textsuperscript{2}-specialized design toward a genuinely abstract subscale
computational-model contract.

\section{Current traceability snapshot}

Table~\ref{tab:stage1-traceability} summarizes the main thesis concepts, their
current computational anchor, and the honest implementation status to be used
in the next writing pass of the thesis.

\begin{table}[htbp]
  \centering
  \caption{Stage-1 traceability snapshot.}
  \label{tab:stage1-traceability}
  \small
  \begin{tabular}{@{}p{0.23\textwidth}p{0.28\textwidth}p{0.12\textwidth}p{0.27\textwidth}@{}}
    \toprule
    \textbf{Thesis concept} & \textbf{Primary code anchor} & \textbf{Status}
    & \textbf{Critical note} \\
    \midrule
    Computational model as method + state + implementation
    & \texttt{Model.hh}, \texttt{NLAnalysis.hh}, \texttt{DynamicAnalysis.hh},
      \texttt{MultiscaleAnalysis.hh}
    & implemented
    & Solver state, restart, and orchestration are first-class citizens. \\

    Multiscale model as aggregation of single-scale models plus coupling laws
    & \texttt{MultiscaleModel.hh}, \texttt{BeamMacroBridge.hh},
      \texttt{MultiscaleTypes.hh}
    & implemented
    & This is currently one of the most thesis-aligned parts of the library. \\

    Material/history separation
    & \texttt{MaterialStatePolicy.hh}, \texttt{Material.hh},
      \texttt{ConstitutiveRelation.hh}
    & implemented
    & Semantically strong, but the runtime-erased name \texttt{Material} still
      overloads the physical meaning of “material”. \\

    Continuum ontology of body/configuration/objective measures
    & \texttt{ContinuumSemantics.hh}, \texttt{ConstitutiveKinematics.hh}, \texttt{KinematicPolicy.hh},
      \texttt{BodyPoint}, \texttt{PlacementMap}, \texttt{MotionSnapshot},
      \texttt{Point.hh}, \texttt{Topology.hh}
    & partial
    & A first compile-time semantic layer now exists, but the geometric and
      topological ontology is still weaker than the continuum chapter. \\

    Total Lagrangian continuum path
    & \texttt{KinematicPolicy.hh}, \texttt{ContinuumElement.hh}
    & implemented
    & This is the current 3D finite-strain reference path. \\

    Updated Lagrangian continuum path
    & \texttt{KinematicPolicy.hh}, \texttt{test\_updated\_lagrangian.cpp}
    & partial
    & Spatial assembly exists, but the validated scientific scope remains
      narrower than the full thesis rhetoric on UL. \\

    Corotational continuum path
    & \texttt{KinematicPolicy.hh}
    & placeholder
    & Exists as a planned slot, not as a usable 3D continuum formulation. \\

    Corotational beam/shell paths
    & \texttt{BeamKinematicPolicy.hh}, \texttt{ShellKinematicPolicy.hh}
    & implemented
    & Must be documented separately from the continuum status. \\

    Abstract local-submodel contract
    & \texttt{SubscaleModelConcepts.hh}, \texttt{LocalModelAdapter.hh}
    & partial
    & Stage-1 now includes a generic contract layer plus a second
      operator-driven realization (\texttt{AffineSectionSubscaleModel.hh}),
      even though the normative orchestrator is still section-specialized. \\
    \bottomrule
  \end{tabular}
\end{table}

\section{Architectural target for Stage 1}

The central architectural move of Stage 1 is to lift the submodel contract by
one abstraction level.  The current production path remains a \emph{section}
FE\textsuperscript{2} realization, but it should no longer be treated as the
only possible shape of a subscale model.

\begin{figure}[htbp]
  \centering
  \begin{tikzpicture}[
      node distance=8mm and 12mm,
      box/.style={draw, rounded corners, align=center, fill=blue!5,
                  minimum width=34mm, minimum height=9mm},
      spe/.style={draw, rounded corners, align=center, fill=green!7,
                  minimum width=42mm, minimum height=10mm},
      impl/.style={draw, rounded corners, align=center, fill=orange!10,
                   minimum width=46mm, minimum height=10mm},
      arr/.style={-{Latex[length=2.5mm]}, thick}
    ]
    \node[box] (step) {StepSolvable\\SubscaleModel};
    \node[box, right=of step] (chk) {Checkpointable\\SubscaleModel};
    \node[box, right=of chk] (obs) {Observable\\SubscaleModel};

    \node[box, below=of chk] (drv) {DrivenSubscaleModel\\$\langle$DrivingState$\rangle$};
    \node[box, right=of drv] (op) {RequestedEffectiveOperatorProvider\\$\langle$Request, Operator$\rangle$};

    \node[spe, below=12mm of drv, xshift=26mm] (sec) {Section-specialized layer\\\texttt{LocalModelAdapter.hh}};
    \node[impl, below=12mm of sec] (nlsm) {Current realization\\\texttt{NonlinearSubModelEvolver}};

    \draw[arr] (step.south) -- (sec.north west);
    \draw[arr] (chk.south) -- (sec.north);
    \draw[arr] (obs.south) -- (sec.north east);
    \draw[arr] (drv.south) -- (sec.north west);
    \draw[arr] (op.south) -- (sec.north east);
    \draw[arr] (sec.south) -- (nlsm.north);
  \end{tikzpicture}
  \caption{Stage-1 contract elevation: a generic subscale-model concept family
  sits above the current section-local FE\textsuperscript{2} specialization.}
  \label{fig:stage1-subscale-contracts}
\end{figure}

The important point is not the new header alone, but the intended
responsibility split:
\begin{itemize}
  \item generic lifecycle and rollback responsibilities are formulation-agnostic;
  \item the driving state and the effective operator become explicit type
  parameters, not hidden assumptions;
  \item the current section FE\textsuperscript{2} route becomes a
  specialization, not the whole ontology;
  \item and type erasure remains an experimental boundary utility rather than
  the normative hot path.
\end{itemize}

\paragraph{First proof that the abstraction is not rhetorical.}
Stage~1 now includes a deliberately simple second realization:
\texttt{AffineSectionSubscaleModel.hh}.  It does not solve a local continuum
problem, does not depend on Ko--Bathe, and does not require a finite-element
submesh.  Instead, it maps a driven generalized state to an affine effective
section operator through a checkpointable, typed lifecycle.  This does not yet
close the broader ambition of supporting XFEM, DG, DPG, or surrogate operators
through one fully generic orchestrator, but it does prove an important point:
the new contract layer is already capable of hosting a non-continuum,
operator-driven local realization without introducing runtime polymorphism into
the hot path.

\section{Reproducible thesis build chain}

The thesis must be treated as a first-class research artifact, not as a
manually compiled side document.  During this cycle, the actual VS Code recipe
used in local development was audited and versioned into the workspace
configuration.  The normative recipe is named \texttt{kaobook} and runs:
\[
\texttt{xelatex}
\rightarrow
\texttt{makeindex}
\rightarrow
\texttt{makeindex\_nomencl}
\rightarrow
\texttt{biber}
\rightarrow
\texttt{makeglossaries}
\rightarrow
\texttt{xelatex}
\rightarrow
\texttt{xelatex}.
\]

This matters for two reasons:
\begin{itemize}
  \item the thesis uses the \texttt{kaobook} class together with
  \texttt{polyglossia}, \texttt{kaobiblio}, glossaries, and nomenclature, so a
  single \texttt{pdflatex} pass is not representative of the real document
  lifecycle;
  \item previous attempts with an incompatible engine or incomplete sequence
  produced misleading failures unrelated to the actual content under review.
\end{itemize}

The versioned workspace recipe does not guarantee that the full manuscript is
warning-free, but it does make the intended compilation chain explicit and
reproducible from the repository itself.

An analogous principle now governs the Windows/MSYS2 regression gate of the
library itself.  The current workflow no longer relies exclusively on
\texttt{PATH}-based runtime discovery for PETSc/VTK-heavy executables.  In
addition to explicit MS-MPI discovery, the CI now stages the full
\texttt{ucrt64/bin/*.dll} runtime surface next to the regression executables.
This is an intentionally conservative operational choice: the audited GitHub
failure mode was \texttt{STATUS\_DLL\_NOT\_FOUND} on the deep
\texttt{fall\_n\_evolver\_advanced\_test.exe} surface, and a broader staging
policy is preferable to repeatedly underestimating transitive runner
dependencies and mistaking an environment failure for a scientific regression.

\section{Finite-kinematics audit to be completed before the FEM rewrite}

Stage 1 must not rewrite the FEM chapter before the finite-kinematics status is
honestly classified.  A purely semantic statement such as ``this formulation
uses a work-conjugate pair'' is not enough: the library also needs an explicit
record of whether the formulation has a real continuum path, whether it has
been regression tested, and whether it should be treated as a reference route
or only as a partial implementation.  This is now encoded in
\texttt{FormulationAuditScope} and exposed through
\texttt{KinematicFormulationTraits<Policy>::audit\_scope}.  Table
\ref{tab:stage1-kinematics-status} records the current target classification.

\begin{table}[htbp]
  \centering
  \caption{Finite-kinematics status to govern the next thesis rewrite.}
  \label{tab:stage1-kinematics-status}
  \small
  \begin{tabular}{@{}llll@{}}
    \toprule
    \textbf{Family} & \textbf{Formulation} & \textbf{Status} & \textbf{Immediate action} \\
    \midrule
    Continuum 3D & Total Lagrangian & Implemented & use as reference baseline \\
    Continuum 3D & Updated Lagrangian & Partial & audit and delimit validated scope \\
    Continuum 3D & Corotational & Placeholder & either implement or explicitly descope \\
    Beams & Corotational & Implemented & preserve and document separately \\
    Shells & Corotational & Implemented & preserve and document separately \\
    \bottomrule
  \end{tabular}
\end{table}

This table matters because the current thesis text already treats TL and UL as
available computational realizations of the continuum theory.  That is broadly
true, but not yet with the same degree of scientific maturity across all
material and element paths.

\section{A fourth audit layer: concrete element--solver pairs}

The previous audit layers remain necessary, but they still leave one
methodological gap: a thesis can honestly state that a formulation is available
for a family and that a solver route exists, while still failing to record
whether the \emph{actual typed pair used in code} is treated as a normative,
partial or merely declared computational path.  Stage~1 now closes part of
that gap through a fourth compile-time layer in
\texttt{src/analysis/ComputationalScopeAudit.hh}.

The intent is deliberately practical.  The library must be able to express,
without runtime overhead, statements such as:
\begin{itemize}
  \item \texttt{ContinuumElement<TL> + NonlinearAnalysis} is the current
  reference finite-kinematics quasi-static path for continuum 3D;
  \item \texttt{ContinuumElement<UL> + NonlinearAnalysis} is real, but not yet
  a reference path;
  \item \texttt{BeamElement<Corotational> + NonlinearAnalysis} has an actual
  runtime route, but still carries a scope disclaimer at the present solver
  evidence level;
  \item \texttt{MITCShellElement<Corotational> + NonlinearAnalysis} is
  available, but not promoted to the same status as the continuum TL baseline.
\end{itemize}

This is a stronger claim than the earlier family matrix because it binds the
audit to the same C++ types that instantiate the numerical kernels.  In other
words, the thesis no longer needs to trust that a family-level table is
implicitly respected by the implementation: the implementation can now expose
that same matrix through the types themselves.

\begin{table}[htbp]
  \centering
  \caption{Representative audited element--solver pairs after the Stage-1
  composition layer.}
  \label{tab:stage1-element-solver-audit}
  \small
  \begin{tabular}{@{}llll@{}}
    \toprule
    \textbf{Typed pair} & \textbf{Scientific reading} & \textbf{Status} &
    \textbf{Interpretation} \\
    \midrule
    \texttt{ContinuumElement<TL>} + \texttt{NonlinearAnalysis}
    & continuum 3D + TL + Newton
    & reference
    & current quasi-static finite-kinematics baseline \\

    \texttt{ContinuumElement<UL>} + \texttt{NonlinearAnalysis}
    & continuum 3D + UL + Newton
    & partial
    & real path with bounded evidence and explicit disclaimer \\

    \texttt{ContinuumElement<small strain>} + \texttt{LinearAnalysis}
    & continuum 3D + linearized + static
    & reference
    & current small-strain baseline \\

    \texttt{BeamElement<corotational>} + \texttt{NonlinearAnalysis}
    & beam 1D + CR + Newton
    & partial
    & runtime path exists, but solver-level promotion remains conservative \\

    \texttt{MITCShellElement<corotational>} + \texttt{NonlinearAnalysis}
    & shell 2D + CR + Newton
    & partial
    & available path, not yet a reference nonlinear shell route \\
    \bottomrule
  \end{tabular}
\end{table}

Conceptually, this fourth layer is important because it is the first one that
lets the software answer a question of the form ``what is the audited status of
the concrete computational model I am actually instantiating right now?''  That
question will become central when Chapter~7 moves from architectural
motivation to a real description of the computational model of
\texttt{fall\_n}.

\section{A fifth audit layer: typed model--solver slices}

The element--solver layer is already much stronger than a purely editorial
matrix, but one final methodological gap remains.  Users do not instantiate
isolated element types in the abstract; they instantiate \emph{model slices}
whose element policy, kinematic formulation, and solver route must remain
coherent.  To make that slice explicit, Stage~1 now adds
\texttt{src/analysis/ComputationalModelSliceAudit.hh}.

This layer composes the previous audit over actual pairs of the form
\[
\texttt{Model<...>} + \texttt{Solver<...>},
\]
so the code can distinguish between
\begin{itemize}
  \item a formulation/route combination that is theoretically admissible,
  \item an element--solver pair that is auditably available, and
  \item the \emph{specific computational slice} that a user is truly running.
\end{itemize}

In practical terms, the new layer can now answer questions such as:
\begin{itemize}
  \item whether \texttt{Model<TL>} + \texttt{NonlinearAnalysis<TL>} is a
  reference finite-kinematics continuum slice;
  \item whether \texttt{Model<UL>} + \texttt{NonlinearAnalysis<UL>} is
  available but still partial;
  \item whether \texttt{Model<TL>} + \texttt{ArcLengthSolver<TL>} should still
  carry a scope disclaimer even though the same model slice is normative under
  Newton continuation.
\end{itemize}

This fifth layer is architecturally modest on purpose.  It does not add a new
runtime registry or a reflective metadata system.  Instead, it reuses
\texttt{element\_type}, \texttt{model\_type}, and existing audited traits as
\texttt{constexpr} compile-time facts.  That keeps the numerical kernels free
from dynamic dispatch while making the thesis language much more precise:
claims can now be anchored not only to ``family/formulation/route'' but to the
same typed computational slices that appear in user code and regression tests.

During the present iteration this family-level audit also stopped being a
purely switch-based internal convention. The library now exposes the same
matrix as structured compile-time data through
\texttt{canonical\_family\_formulation\_audit\_table()},
\texttt{canonical\_family\_formulation\_audit\_row(...)} and the
\texttt{find\_family\_*reference\_path(...)} helpers. This matters for the
thesis because the FEM chapter can now cite a single computational source of
truth rather than reproduce the family/formulation classification manually in
multiple places.

\paragraph{Family-level scope matrix.}
The previous semantic lift was formulation-centric.  Stage~1 now adds a second,
strictly separate layer in \texttt{FormulationScopeAudit.hh}, because the same
formulation label does not mean the same thing across all element families.  In
particular, \emph{corotational} is still a placeholder for continuum 3D, while
beam corotational is the current reference nonlinear beam path and shell
corotational is available but still not promoted to the same baseline status.
This matrix is no longer merely documentary: the library now exposes
\texttt{FamilyKinematicPolicyAuditTraits} together with the concepts
\texttt{FamilyNormativelySupportedKinematicPolicy} and
\texttt{FamilyReferenceGeometricNonlinearityKinematicPolicy}, so the audited
scope can also act as a compile-time compatibility layer.

\begin{table}[htbp]
  \centering
  \caption{Family-level formulation audit encoded in
  \texttt{canonical\_family\_formulation\_audit\_scope(...) }.}
  \label{tab:stage1-family-formulation-audit}
  \small
  \begin{tabular}{@{}llll@{}}
    \toprule
    \textbf{Family} & \textbf{Formulation} & \textbf{Current code status}
    & \textbf{Critical reading} \\
    \midrule
    Continuum 3D & Small strain & implemented & linear reference path only \\
    Continuum 3D & Total Lagrangian & reference baseline & current nonlinear reference \\
    Continuum 3D & Updated Lagrangian & partial & real path, still under scope disclaimer \\
    Continuum 3D & Corotational & placeholder & not a usable 3D continuum formulation \\
    Beam 1D & Small rotation & implemented & linearized beam baseline \\
    Beam 1D & Corotational & reference baseline & current nonlinear beam reference \\
    Shell 2D & Small rotation & reference baseline & current structural-shell baseline \\
    Shell 2D & Corotational & partial & available and tested, but not yet promoted to the same status \\
    \bottomrule
  \end{tabular}
\end{table}

This family-level matrix is no longer backed only by concept-level harnesses.
During the present iteration it was also rechecked on structural targets that
instantiate the normative templates directly: the beam regression surface
(\texttt{fall\_n\_beam\_test.exe}, \(21/21\)), the MITC-shell regression
surface (\texttt{fall\_n\_mitc\_shell\_test.exe}, \(12/12\)), and the seismic
infrastructure harness (\texttt{fall\_n\_seismic\_infra\_test.exe},
\(33/33\)) remained in green after the family-aware constraints were moved
into \texttt{ContinuumElement}, \texttt{BeamElement}, and
\texttt{MITCShellElement}.  Methodologically, this matters because it shows
that the matrix is already acting as a compile-time boundary on the real
structural types, not only as a declarative table attached to the thesis.

\paragraph{A third audit layer: analysis routes.}
The family--formulation matrix is necessary, but it is still not sufficient.
A formulation being available for a family does not imply that every solver
route is equally mature for that same scope.  The current codebase already
distinguishes at least four routes with different lifecycle and evidence
profiles:
\begin{itemize}
  \item \texttt{LinearAnalysis.hh} for linear-static solves;
  \item \texttt{NLAnalysis.hh} for checkpointable incremental Newton solves;
  \item \texttt{DynamicAnalysis.hh} for implicit second-order dynamics;
  \item \texttt{ArcLengthSolver.hh} for continuation through limit points.
\end{itemize}

Stage~1 now records that distinction explicitly in
\texttt{src/analysis/AnalysisRouteAudit.hh}.  This layer classifies
\emph{family \(\times\) formulation \(\times\) route} rather than only
\emph{family \(\times\) formulation}, and it records whether the route has a
runtime path, whether it supports checkpoint/restart or directed stepping, and
whether the current evidence is strong enough to treat it as a reference
baseline.

\begin{table}[htbp]
  \centering
  \caption{Representative audited scopes after adding the analysis-route layer.}
  \label{tab:stage1-family-formulation-route-audit}
  \small
  \begin{tabular}{@{}lllll@{}}
    \toprule
    \textbf{Family} & \textbf{Formulation} & \textbf{Route} & \textbf{Status}
    & \textbf{Critical reading} \\
    \midrule
    Continuum 3D & Small strain & Linear static & reference baseline
    & audited linear reference route \\
    Continuum 3D & Total Lagrangian & Nonlinear Newton & reference baseline
    & current finite-kinematics reference path \\
    Continuum 3D & Updated Lagrangian & Nonlinear Newton & partial
    & real route under validated-scope disclaimer \\
    Continuum 3D & Total Lagrangian & Implicit dynamics & interface declared
    & generic route exists, but not yet a reference claim \\
    Continuum 3D & Total Lagrangian & Arc-length & partial
    & implemented continuation class, still not normative baseline \\
    Beam 1D & Small rotation & Nonlinear Newton & implemented
    & structural route already exercised in current tests \\
    Beam 1D & Corotational & Nonlinear Newton & partial
    & formulation mature, solver-level evidence still lagging \\
    Shell 2D & Corotational & Nonlinear Newton & partial
    & element path exists, solver route not yet baseline \\
    \bottomrule
  \end{tabular}
\end{table}

This third layer produces a useful scientific correction: it is still true
that beam and shell corotational kinematics exist and matter, but it is no
longer acceptable to infer from that fact alone that every global nonlinear
solver route using those formulations has reached the same maturity as the
continuum Total-Lagrangian/Newton baseline.

\section{WP3 initial delivery: a first semantic lift of continuum theory}

Stage~1 now includes a first explicit semantic layer in
\texttt{src/continuum/ContinuumSemantics.hh}.  The purpose of this layer is not
to introduce a new runtime object model, but to encode three distinctions that
were already present in Chapter~4 and only implicit in code:
\begin{enumerate}
  \item the \emph{description kind} (\emph{linearized}, \emph{material},
  \emph{spatial}, or \emph{corotated});
  \item the configuration over which a formulation is assembled; and
  \item the strain--stress pair that is being treated as work-conjugate.
\end{enumerate}

This layer now goes one step further by distinguishing \emph{mathematical
semantics} from \emph{audited computational scope}.  In practice, the code now
stores:
\begin{itemize}
  \item \texttt{VirtualWorkSemantics}, which captures the intended virtual-work
  statement of a formulation;
  \item \texttt{AuditEvidenceLevel}, which classifies the strength of the
  current evidence (\emph{interface declared}, \emph{regression tested},
  \emph{reference baseline}, etc.); and
  \item \texttt{FormulationAuditScope}, which records whether the formulation
  actually has a continuum 3D path, whether it supports finite kinematics, and
  whether it should be used as the current reference route.
\end{itemize}

The practical consequence is important: the library can now state, at compile
time, that Total Lagrangian is the current reference path for finite-kinematic
continuum work, that Updated Lagrangian is real but still partial, and that
continuum corotational remains a placeholder without pretending that all three
carry the same computational weight.

Mathematically, this corresponds to making explicit the different virtual-work
pairs already used throughout the library:
\begin{align}
  \delta W_{\mathrm{lin}} &= \int_{\Omega_0} \boldsymbol{\sigma} :
    \delta\boldsymbol{\varepsilon}\,\mathrm{d}V,\\
  \delta W_{\mathrm{TL}} &= \int_{\Omega_0} \mathbf{S} :
    \delta\mathbf{E}\,\mathrm{d}V,\\
  \delta W_{\mathrm{UL}} &= \int_{\Omega_t} \boldsymbol{\sigma} :
    \delta\mathbf{e}\,\mathrm{d}v,
\end{align}
where \(\boldsymbol{\varepsilon}\) is the infinitesimal strain,
\(\mathbf{E}\) the Green--Lagrange strain, \(\mathbf{e}\) the Almansi strain,
\(\mathbf{S}\) the second Piola--Kirchhoff stress, and
\(\boldsymbol{\sigma}\) the Cauchy stress.

The key architectural point is that the code now records these distinctions as
compile-time semantics and not only as comments or implicit conventions:
\begin{itemize}
  \item \texttt{KinematicFormulationTraits<Policy>} classifies each formulation
  by description kind, assembly configuration, canonical conjugate pair, and
  scientific maturity;
  \item \texttt{BodyPoint}, \texttt{ReferencePoint}, \texttt{CurrentPoint},
  \texttt{PlacementMap} and \texttt{MotionSnapshot} provide a first
  computational vocabulary for the Chapter~4 notions of body, configuration,
  placement, and motion;
  \item \texttt{VirtualWorkSemantics} and its associated compatibility/mass-
  measure tags make explicit whether the formulation claims an exact
  work-conjugate statement, only a linearized-equivalent one, or an unaudited
  placeholder status;
  \item \texttt{ConstitutiveKinematics<dim>} now carries the corresponding
  semantic metadata together with the actual strain tensors, \(F\), and \(J\);
  \item tests verify both numerical consistency and semantic consistency; and
  \item the element templates \texttt{ContinuumElement},
  \texttt{BeamElement}, and \texttt{MITCShellElement} now consume the family
  audit layer directly, so unsupported formulation/family combinations are
  rejected at instantiation time instead of being left to narrative warnings.
\end{itemize}

This is intentionally modest.  The semantic lift does \emph{not} yet close the
entire Chapter~4 ontology in software.  It does, however, establish a truthful
bridge between continuum theory and implementation, which is a necessary
precondition for the later rewrite of the FEM chapter.

\paragraph{Why the new carriers matter.}
The added point/placement/motion carriers are intentionally shallow. They do
not try to replace the existing geometric kernel, nor do they claim to encode
the full differential-geometric structure of the continuum chapter. Their
purpose is narrower and more pragmatic:
\begin{itemize}
  \item to distinguish, in code, between a point on the material body and its
  images in reference/current/corotated configurations;
  \item to make the placement gradient and its Jacobian explicit when a local
  continuum state is passed through the constitutive interface; and
  \item to prepare future compile-time checks on work-conjugacy, objectivity,
  and formulation compatibility without adding polymorphic cost to the hot
  path.
\end{itemize}
This is exactly the sort of semantic lift that the thesis needs before the FEM
chapter can be rewritten in terms of the library itself instead of in generic
textbook language.

\paragraph{A subtle but important refinement.}
The new virtual-work metadata is not redundant with the strain/stress pair
alone.  Small strain, Total Lagrangian, and Updated Lagrangian all rely on
work-conjugate pairs in a broad mechanical sense, but they do \emph{not} make
the same claim:
\begin{itemize}
  \item TL and UL are treated as \emph{exact} statements on their natural
  domains \(\Omega_0\) and \(\Omega_t\), respectively;
  \item small strain is treated as a \emph{linearized-equivalent} statement,
  because the use of \(\boldsymbol{\sigma}\) with
  \(\boldsymbol{\varepsilon}\) rests on the kinematic simplification
  \(\Omega_t \approx \Omega_0\);
  \item continuum corotational remains an \emph{unaudited placeholder} rather
  than a fully claimed virtual-work formulation.
\end{itemize}
This distinction is mathematically small but architecturally very important:
it prevents the library from flattening all formulations into the same
semantic bucket just because they share parts of the implementation pipeline.

\section{WP4 interim audit: what the formulation status really means}

The formulation labels now have an explicit computational meaning:
\begin{itemize}
  \item \textbf{implemented} means the formulation exists in code, participates
  in the normative path, and has dedicated tests supporting its current claim;
  \item \textbf{partial} means the formulation exists and is useful, but its
  validated scientific scope is narrower than a full theoretical treatment;
  \item \textbf{placeholder} means the architecture reserves the slot, but the
  library does not yet provide a mature scientific realization.
\end{itemize}

For the present codebase, this leads to the following interpretation:
\begin{itemize}
  \item \textbf{Total Lagrangian continuum 3D} remains the finite-strain
  reference path.
  \item \textbf{Updated Lagrangian continuum 3D} is no longer called a
  placeholder.  The spatial pathway, push-forward machinery, and dedicated
  tests are real.  However, it remains \emph{partial} because the validated
  scope is still narrower than the full rhetoric of a general UL continuum
  formulation.
  \item \textbf{Continuum corotational} remains a placeholder.  The beam and
  shell corotational implementations must not be used as indirect evidence that
  continuum corotational kinematics is already available.
\end{itemize}

This distinction is important for the thesis: it prevents the FEM chapter from
making scientifically stronger claims than the current evidence supports.

\section{Stage-1 work packages}

\subsection*{WP1 -- Formal traceability audit}
\textbf{Goal:} create a chapter-by-chapter map from thesis concepts to code,
tests, and evidence. \\
\textbf{Outputs:} status tables, honest scope statements, and audit notes in
both the thesis and \texttt{fall\_n/doc}. \\
\textbf{Exit criterion:} no central thesis claim remains uncategorized as
implemented, partial, placeholder, or future work.

\subsection*{WP2 -- Contract elevation for local submodels}
\textbf{Goal:} introduce a generic compile-time contract family above the
current section-local FE\textsuperscript{2} adapter. \\
\textbf{Outputs:} generic concepts, section-specialized adapters, and a staged
migration plan for the multiscale orchestrator. \\
\textbf{Exit criterion:} the codebase can host more than one kind of local
subproblem without pushing virtual dispatch into the normative hot path.

\subsection*{WP3 -- Continuum-mechanics semantic lift}
\textbf{Goal:} encode more faithfully the ontology of body, configuration,
placement, and conjugate pairs from the continuum chapter. \\
\textbf{Outputs:} tags, carriers, and semantic bridges from theory to
implementation. \\
\textbf{Exit criterion:} the finite-kinematics pipeline can be explained and
tested in the same language as the continuum chapter.

\subsection*{WP4 -- Finite-kinematics verification campaign}
\textbf{Goal:} revalidate TL, UL, and corotational paths with formulation-level
honesty. \\
\textbf{Outputs:} dedicated tests, explicit scope tables, and formulation
status records. \\
\textbf{Exit criterion:} the FEM chapter can be rewritten without overstating
UL or continuum corotational support.

\subsection*{WP5 -- FEM chapter rewrite}
\textbf{Goal:} rewrite Chapter 6 as the theory of the library rather than a
generic FEM survey. \\
\textbf{Outputs:} a chapter centered on domain topology, kinematic policies,
constitutive integration, assembly, tangent operators, and solver architecture. \\
\textbf{Exit criterion:} the thesis and the library speak the same FEM
language.

\subsection*{WP6 -- Implementation chapter write-up}
\textbf{Goal:} turn the current thesis implementation chapter into a true
description of the computational model. \\
\textbf{Outputs:} architecture diagrams, concept families, lifecycle logic,
HPC remarks, and multiscale orchestration rationale. \\
\textbf{Exit criterion:} Chapter 7 can stand as a rigorous software-architecture
chapter for a scientific-computing thesis.

\subsection*{WP7 -- Thesis closure preparation}
\textbf{Goal:} define the remaining validation and conclusion chapters. \\
\textbf{Outputs:} a projected Chapter 8 (verification/validation), Chapter 9
(cyclic multiscale validation and computational performance), and Chapter 10
(conclusions and outlook). \\
\textbf{Exit criterion:} the thesis can transition from architectural
consolidation to final scientific validation.

\section{Critical publication outlook}

At this point, the material is \emph{promising} for high-level publications,
but not yet automatically deserving of them.

\begin{itemize}
  \item A strong methodology paper becomes realistic only after the
  multiscale validation closes the post-cracking regime convincingly.
  \item A software/architecture paper is already plausible, but it should be
  marketed honestly as an evolving scientific-computing architecture, not as a
  fully generalized multiscale platform.
  \item A paper claiming a general abstract local-model framework would still be
  premature until a second genuinely different local-model family passes
  through the elevated contract.
\end{itemize}

In other words, the present stage justifies a roadmap and architectural
hardening effort; it does not yet justify complacency about publication
readiness.
